% Standalone problem document, generated from the adjacent README.md.
% Compile with XeLaTeX. Mathematical content is maintained in Markdown.
\documentclass[11pt,a4paper]{article}
\usepackage[fontsize=10.5pt]{scrextend}
\usepackage[margin=22mm,headheight=15pt,headsep=9mm,footskip=11mm]{geometry}
\usepackage{fontspec}
\setmainfont{texgyrepagella}[Extension=.otf,UprightFont=*-regular,BoldFont=*-bold,ItalicFont=*-italic,BoldItalicFont=*-bolditalic]
\setsansfont{texgyreheros}[Extension=.otf,UprightFont=*-regular,BoldFont=*-bold,ItalicFont=*-italic,BoldItalicFont=*-bolditalic]
\setmonofont{texgyrecursor}[Extension=.otf,UprightFont=*-regular,BoldFont=*-bold,ItalicFont=*-italic,BoldItalicFont=*-bolditalic,Scale=MatchLowercase]
\usepackage{amsmath,amssymb}
\usepackage{unicode-math}
\setmathfont{texgyrepagella-math.otf}
\usepackage{xcolor,microtype,enumitem,needspace,fancyhdr}
\usepackage{bookmark}
\definecolor{ink}{HTML}{17344B}
\definecolor{accent}{HTML}{197D80}
\definecolor{muted}{HTML}{596773}
\hypersetup{colorlinks=true,linkcolor=accent,urlcolor=accent,pdftitle={SP-07 - The
sharp spectral-matching constant for normal
matrices},pdfauthor={Open Problems in Numerical Linear Algebra}}
\urlstyle{same}
\setlength{\parindent}{0pt}
\setlength{\parskip}{5pt plus 1pt minus 1pt}
\linespread{1.04}
\setlength{\emergencystretch}{3em}
\widowpenalty=10000
\clubpenalty=10000
\displaywidowpenalty=10000
\allowdisplaybreaks[1]
\setlist{nosep,leftmargin=1.5em}
\providecommand{\tightlist}{\setlength{\itemsep}{0pt}\setlength{\parskip}{0pt}}
\providecommand{\pandocbounded}[1]{#1}
\pagestyle{fancy}
\fancyhf{}
\fancyhead[L]{\sffamily\footnotesize\color{muted}OPEN PROBLEMS IN NUMERICAL LINEAR ALGEBRA}
\fancyhead[R]{\sffamily\bfseries\footnotesize\color{accent}SP-07}
\fancyfoot[L]{\parbox[b]{0.9\textwidth}{\sffamily\fontsize{8}{10}\selectfont\color{muted}Editorial ratings; a bounded search cannot certify that a problem remains unsolved.\\Literature check: 2026-09-10}}
\fancyfoot[R]{\sffamily\footnotesize\color{muted}\thepage}
\renewcommand{\headrulewidth}{0.3pt}
\renewcommand{\headrule}{\hbox to\headwidth{\color{accent}\leaders\hrule height \headrulewidth\hfill}}
\makeatletter
\renewcommand{\section}{\Needspace{5\baselineskip}\@startsection{section}{1}{0pt}{12pt plus 2pt minus 2pt}{4pt}{\sffamily\bfseries\large\color{ink}}}
\renewcommand{\subsection}{\Needspace{4\baselineskip}\@startsection{subsection}{2}{0pt}{9pt plus 1pt minus 1pt}{3pt}{\sffamily\bfseries\normalsize\color{ink}}}
\makeatother
\setcounter{secnumdepth}{0}
\begin{document}
{\sffamily\small\color{muted}Eigenvalues and inverse problems\par}
\vspace{3pt}
{\raggedright\hyphenpenalty=10000\exhyphenpenalty=10000\sffamily\bfseries\fontsize{21}{24}\selectfont\color{ink}The
sharp spectral-matching constant for normal matrices\par}
\vspace{7pt}
{\sffamily\small\textbf{Difficulty:} extreme\quad\textbf{Importance:} interesting
to the community\par}
{\sffamily\small\bfseries\color{accent}Status: Open\par}
\vspace{4pt}
{\color{accent}\hrule height 0.8pt}
\vspace{6pt}
\textbf{Rating rationale:} Extreme reflects a longstanding
sharp-constant problem beyond the false constant-one conjecture;
community impact is eigenvalue matching and perturbation theory for
normal matrices.

\subsection{Problem statement}\label{problem-statement}

For normal matrices \(A,B\in\mathbb C^{n\times n}\), list their
eigenvalues \(\alpha_1,\ldots,\alpha_n\) and \(\beta_1,\ldots,\beta_n\)
with algebraic multiplicity. Define the bottleneck matching distance

\[
d_\infty(A,B)=\min_{\pi\in S_n}\max_{1\le j\le n}
|\alpha_j-\beta_{\pi(j)}|,
\]

where \(S_n\) is the permutation group. Determine the exact value of

\[
C_{\rm normal}=\sup_{n\ge2}\ \sup_{\substack{A,B\in\mathbb C^{n\times n}\ \mathrm{normal}\\ A\ne B}}
\frac{d_\infty(A,B)}{\|A-B\|_2}.
\]

Here normal means \(AA^*=A^*A\), and \(\|\cdot\|_2\) is the spectral
norm. The target is a single sharp constant valid in every dimension.
Equivalently, find the smallest \(C\) for which
\(d_\infty(A,B)\le C\|A-B\|_2\) for all such pairs, with matching
sharpness evidence. The conjectured value \(C=1\) has already been
refuted; it is not the question being posed.

\subsection{Why it matters}\label{why-it-matters}

This is a sharp conditioning question for the complete spectrum of a
normal matrix. Matching counts repeated eigenvalues, so a
Hausdorff-distance estimate that ignores multiplicity does not answer
it.

\newpage
\subsection{References}\label{references}

R. Bhatia,
\href{https://epubs.siam.org/doi/10.1137/1.9780898719079.ch9}{\emph{Perturbation
Bounds for Matrix Eigenvalues}}, SIAM, 2007, Chapter IX, pp.~154--155;
the question is restated as Problem 27 in X. Zhan,
\href{https://math.ecnu.edu.cn/~zhan/papers/ZhanICCM.pdf}{\emph{Open
problems in matrix theory}}, §21, pp.~12--13. P. Šemrl,
\href{https://indico.ictp.it/event/a08167/session/73/contribution/52/material/0/0.pdf}{\emph{Geometrical
techniques in linear algebra}}, ICTP lecture notes, 2009, p.~10, Theorem
5 and the following open question. J. A. Holbrook,
\href{https://doi.org/10.1016/0024-3795(92)90047-E}{\emph{Spectral
variation of normal matrices}}, LAA 174 (1992), 131--144.

\subsection{Earlier status check ---
2026-09-08}\label{earlier-status-check-2026-09-08}

The cited sources distinguish the unresolved sharp constant from
Holbrook's counterexample to one. They give \(1<C_{\rm normal}<3\). A.
Parusiński and A. Rainer,
\href{https://arxiv.org/abs/2603.23056}{\emph{Eigenvalue stability of
Hermitian and normal matrices}}, v1 (2026), Proposition 3.4, still
states the classical universal bound with \(1<C<3\); it does not
determine the sharp constant. Searches for
\textit{sharp spectral variation constant normal matrices},
\textit{spectral variation best possible constant}, and 2025/2026
follow-ups found no exact value. The explicit open-question sources are
historical; the recent bound is context, not an explicit reaffirmation
of openness. This is a bounded literature check.

\subsection{Audit update --- 2026-09-10}\label{audit-update-2026-09-10}

The \href{https://arxiv.org/pdf/2603.23056}{2026 spectral perturbation
manuscript}, Proposition 3.4, still uses a universal constant strictly
between 1 and 3 rather than determining the optimum.
Sharp-normal-spectral-matching searches found no exact value. This
supports the retained open verdict alongside the historical lower-bound
examples cited above; the old constant-one statement is already false
and is not the target here.
\end{document}
